\uuid{DkBN}
\titre{Application linéaire/Nombres complexes}
\auteur{hanani ; mzari}
\organisation{exo7}
\licence{CC BY-NC-SA 4.0}
\datecreate{2020-01-20}
\niveau{L1}
\module{Algèbre}
\chapitre{Application linéaire}
\sousChapitre{Définition}
\difficulte{2}
\theme{Applications linéaires, Application linéaire, Nombres complexes}
\exo7qcmid{lille-2.applications-lineaires.006}

\contenu{
\begin{question}
\qclassification{107.01}
\qtags{motcle=Application linéaire/Nombres complexes}
On considère les applications suivantes :
$$\begin{array}{rccc}f:&\Cc&\to&\Cc\\
& z&\to& |z|\end{array}  \quad \mbox{et} \quad \begin{array}{rccc}g:&\Cc&\to&\Cc\\
& z&\to& \overline{z},
\end{array} $$
où $|z|$ (resp. $\overline{z}$) est le module (resp. le conjugué) de $z$. Quelles sont les assertions vraies ?
\begin{answers}
\bad{$f$ est $\Cc$-linéaire}.
\bad{$f$ est $\Rr$-linéaire}.
\good{$g$ est $\Rr$-linéaire}.
\bad{$g$ est $\Cc$-linéaire}.
\end{answers}
\begin{explanations} On vérifie que $f$ n'est pas $\Rr$-linéaire (donc n'est pas $\Cc$-linéaire) et que $g$ est $\Rr$-linéaire, mais non $\Cc$-linéaire.
\end{explanations}
\end{question}
}
